\begin{python0}
from solutions import *; clear()
\end{python0}
\sectionthree{\texttt{switch}}

Now I' ll show you something that' s very
similar to the if and if-else statement. Watch out! There are
\EMPHASIZE{subtle differences} ....

OK. First run this:
\begin{console}
int x = 0;
std::cin >> x;
if (x == 1)
{
     std::cout << "yo" << std::endl;
}
else if (x == 2)
{
     std::cout << "right ..." << std::endl;
}
else if (x == 3)
{
     std::cout << "oh?" << std::endl;
}
else
{
     std::cout << "huh?" << std::endl;
}
std::cout << "out of if-else" << std::endl;
\end{console}

Run this with x = 0, x = 1, x = 2, x = 3, x = 4. No surprises right?

Note that all the conditions uses the == boolean operator and it is used
to compare integer values.

Modify the above:

\begin{console}[commandchars=\~\@\$]
int x = 0;
std::cin >> x;
~textbf@switch (x)$
~textbf@{$
     ~textbf@case 1:$
        std::cout << "yo" << std::endl;
        ~textbf@break;$

     ~textbf@case 2:$
        std::cout << "right ..." << std::endl;
        ~textbf@break;$
        
     ~textbf@case 3:$
        std::cout << "oh?" << std::endl;
        ~textbf@break;$
        
     ~textbf@default:$
        std::cout << "huh?" << std::endl;
~textbf@}$
std::cout << "out of switch" << std::endl;
\end{console}


Run it again.

The \texttt{switch} statement is like a switchboard, directing the flow of
execution based on the integer value of a variable. The flow of
execution is passed to the case whose value matches the value of the
variable.

Now modify it again by getting rid of the \texttt{break} statement in case 2:
\begin{console}
int x = 0;
std::cin >> x;
switch (x)
{
     case 1:
        std::cout << "yo" << std::endl;
        break;

     case 2:
        std::cout << "right ..." << std::endl;
        
     case 3:
        std::cout << "oh?" << std::endl;
        break;
        
     default:
        std::cout << "huh?" << std::endl;
}
std::cout << "out of switch" << std::endl;
\end{console}


What does the \texttt{break} statement do? If you do \EMPHASIZE{not} have a
\texttt{break} statement, execution will actually \EMPHASIZE{go to the statements of the next case without checking the value for that case!!!}

There are \EMPHASIZE{two} ways to get \EMPHASIZE{out} of a \texttt{switch}:
either through a \texttt{\textbf{break}} statement or when execution has
reached the \EMPHASIZE{end of the \texttt{switch} block}.

By now, I don' t have to tell you that you can put any
statement in the switch statement. You can put an if or an if-else into
the switch. You can put a switch into a switch. You can put a switch
into an if-else, etc. Isn' t life ever so colorful?

